% Part: first-order-logic
% Chapter: sequent-calculus
% Section: rules-and-proofs

\documentclass[../../../include/open-logic-section]{subfiles}

\begin{document}

\iftag{FOL}
      {\olfileid{fol}{seq}{rul}}
      {\olfileid{pl}{seq}{rul}}

\olsection{నియమాలు, \usetoken{P}{derivation}}

కింది చర్చలో $\Gamma, \Delta, \Pi, \Lambda$లు \tetoken{వాక్యాలు}{!!{sentence}s} యొక్క
పరిమిత క్రమాలను సూచిస్తాయని అనుకుందాం.

\begin{defn}[సీక్వెంట్]
\emph{సీక్వెంట్} అంటే కింది రూపంలోని వ్యక్తీకరణ:
\[
\Gamma \Sequent \Delta
\]
ఇక్కడ $\Gamma$, $\Delta$లు $\Lang L$ భాషలోని \tetoken{వాక్యాలు}{!!{sentence}s} యొక్క
పరిమిత (శూన్యమైనవీ కావచ్చు) క్రమాలు. $\Gamma$ను \emph{పూర్వాంగం}
అనీ, $\Delta$ను \emph{ఉత్తరాంగం} అనీ అంటారు.
\end{defn}

\begin{explain}
సీక్వెంట్ వెనుకనున్న అంతర్బుద్ధి ఇది: పూర్వాంగంలోని \tetoken{వాక్యాలు}{!!{sentence}s}
అన్నీ చెల్లితే, ఉత్తరాంగంలోని \tetoken{వాక్యాలలో}{!!{sentence}s} కనీసం ఒకటి చెల్లుతుంది.
అంటే, $\Gamma = \tuple{!A_1, \dots, !A_m}$ మరియు
$\Delta = \tuple{!B_1, \dots, !B_n}$ అయితే, కింది వాక్యం
చెల్లినప్పుడు, అప్పుడు మాత్రమే $\Gamma \Sequent \Delta$ చెల్లుతుంది:
\[
(!A_1 \land \cdots \land !A_m) \lif (!B_1 \lor \cdots \lor
!B_n)
\]
రెండు ప్రత్యేక సందర్భాలు ఉన్నాయి: $\Gamma$ శూన్యంగా ఉండడం,
$\Delta$ శూన్యంగా ఉండడం. $\Gamma$ శూన్యంగా, అంటే $m = 0$గా ఉంటే,
$!B_1 \lor \dots \lor !B_n$ చెల్లినప్పుడు, అప్పుడు మాత్రమే $\quad
\Sequent \Delta$ చెల్లుతుంది. $\Delta$ శూన్యంగా, అంటే $n = 0$గా
ఉంటే, $\lnot(!A_1 \land \dots \land !A_m)$ చెల్లినప్పుడు, అప్పుడు
మాత్రమే $\Gamma \Sequent \quad$ చెల్లుతుంది. దానికి సరిపోలే
\tetoken{వాక్యం}{!!{sentence}} చెల్లుబాటైతే, అప్పుడు మాత్రమే సీక్వెంట్ చెల్లుబాటని అంటాం.
\end{explain}

$\Gamma$ ఒక \tetoken{వాక్యాలు}{!!{sentence}s} క్రమమైతే, $\Gamma$ కుడి చివర $!A$ను
చేర్చగా వచ్చే క్రమాన్ని $\Gamma, !A$ అనీ ($\Gamma$ ఎడమ చివర $!A$ను
చేర్చగా వచ్చే క్రమాన్ని $!A, \Gamma$ అనీ) రాస్తాం. $\Delta$ కూడా ఒక
\tetoken{వాక్యాలు}{!!{sentence}s} క్రమమైతే, ఈ రెండు క్రమాల సంయోజనమే $\Gamma, \Delta$.

\begin{defn}[ప్రారంభ సీక్వెంట్]
\emph{ప్రారంభ సీక్వెంట్} అంటే
\iftag{prvFalse,prvTrue}{కింది రూపాల్లో ఒకదానిలోని సీక్వెంట్:
  \begin{enumerate}
    \item $!A \Sequent !A$
    \tagitem{prvTrue}{$\quad \Sequent \ltrue$}{}
    \tagitem{prvFalse}{$\lfalse \Sequent \quad$}{}
  \end{enumerate}}
{భాషలోని ఏ \tetoken{వాక్యం}{!!{sentence}} $!A$కైనా $!A \Sequent !A$ రూపంలోని
సీక్వెంట్}.
\end{defn}

సీక్వెంట్ కలనంలోని \tetoken{వ్యుత్పత్తులు}{!!^{derivation}s} కొన్ని సీక్వెంట్ వృక్షాలు.
వాటిలో అన్నిటికంటే పైనున్నవి ప్రారంభ సీక్వెంట్లు; ఒక సీక్వెంట్ ఒకటి
లేదా రెండు ఇతర సీక్వెంట్ల కింద ఉంటే, అది నిగమన నియమాన్ని సరిగ్గా
వర్తింపజేయడం ద్వారా వాటి నుంచి రావాలి. $\Log{LK}$ నియమాలు రెండు ప్రధాన
రకాలుగా విభజించబడతాయి: \emph{తార్కిక} నియమాలు, \emph{నిర్మాణాత్మక}
నియమాలు. కింది సీక్వెంట్‌లో $!A$ను మరియు/లేదా $!B$ను కలిగిన
\tetoken{వాక్యం}{!!{sentence}} యొక్క \tetoken{ప్రధాన సంయోజకం}{!!{main
  operator}} పేరును తార్కిక నియమానికి పెడతాం.
ప్రతి నియమానికి రెండు రూపాలు ఉంటాయి: \tetoken{తార్కిక సంయోజకాన్ని}{!!{operator}} కలిగిన \tetoken{వాక్యం}{!!{sentence}}
ఎడమవైపు ఉండే సీక్వెంట్‌ను నిగమించే రూపం ఒకటి, అది కుడివైపు ఉండే
\tetoken{వాక్యం}{!!{sentence}} గల సీక్వెంట్‌ను నిగమించే రూపం మరొకటి.

\end{document}
